\begin{qfdefinition}[Private equity: features, stages and strategies, how funds operate, and its investment characteristics]\label{def:private-equity-features}
\qffield{Statement} \emph{Features.} Private equity buys stakes in companies that are not listed, or delists them, through funds that are limited partnerships with a life of about ten years \citep[pp.~594--595]{s7-ang}. Four characteristics distinguish it \citep[pp.~5--7]{metrick-yasuda-survey-2011}:
\begin{enumerate}
\item the managers are financial intermediaries investing other people's money;
\item they invest in private companies, about which public information is scarce, so holdings are not marked to market;
\item they monitor and help manage the companies, through board seats, veto rights and contingent control rights;
\item they must exit, by an IPO, a sale to another investor or a sale to a company, to return cash within the fund's life.
\end{enumerate}

\emph{Strategies} (Ang, pp.~595--596): \emph{venture capital} takes minority stakes in young companies with little or no leverage; \emph{buyouts} acquire control of mature companies with high leverage, and represent most of the money; \emph{growth} capital takes minority stakes in expanding profitable companies; distressed and turnaround funds buy troubled companies; secondary funds buy existing fund stakes; LPs also co-invest.

\emph{Value creation in buyouts.} Kaplan and Str\"omberg group the changes private-equity firms make as \emph{financial engineering} (management equity and leverage, which disciplines the use of free cash flow), \emph{governance engineering} (active, smaller boards) and \emph{operational engineering} (industry and operating expertise) \citep[NBER WP 14207, pp.~13--16]{pe-kaplan-stromberg-2009}. Surveyed buyout managers rank growth in the business and operational improvements above leverage and multiple arbitrage as return drivers \citep[NBER WP 21133, p.~27]{w5-gompers-kaplan-mukharlyamov}.

\emph{Investment characteristics.}
\begin{itemize}
\item \emph{Illiquidity and the J-curve}: capital is called over years and returned on exits; early fees and investments at cost make interim returns negative \citep[pp.~16--18]{ippolito-private-capital-2020}.
\item \emph{Returns and dispersion}: in Burgiss data on nearly 1,400 US funds, average buyout funds beat the S\&P~500 by 20--27 percent over a fund's life, more than 3 percent a year; venture funds outperformed in the 1990s and underperformed in the 2000s \citep[abstract]{pe-harris-jenkinson-kaplan-2014}. Returns vary widely across funds and persist from one fund of a manager to the next \citep[NBER WP 9807, pp.~1--2]{pe-kaplan-schoar-2005}.
\item \emph{Cyclicality}: commitments and buyout activity rise in booms, when expected returns are lowest \citep[p.~596]{s7-ang}; \citep[p.~3]{pe-kaplan-stromberg-2009}.
\item \emph{Holding period}: the median buyout is held about six years; 38 percent of exits are sales to strategic buyers, 24 percent secondary buyouts, 14 percent IPOs (pp.~11--12).
\end{itemize}

\qffield{Practice} \emph{Ashby Fund I} (original, \hyperref[meth:private-fund-multiples-pme]{metrics}): four buyouts bought in years 0--3 are worth 0.93 of paid-in capital after two years, because of fees and holding at cost. Exits in years 4--7 lift the investor's multiple to 1.52, with the since-inception IRR rising from $-14.8$ percent to 11.7 percent.

\qffield{Conditions} The return evidence covers US funds over specific vintages and is net of fees but not adjusted for leverage or market beta beyond the S\&P~500 comparison. The NBER working-paper page numbers are used for the papers marked so.

\qffield{Pitfalls} Treating average performance as attainable when dispersion and access dominate outcomes. Reading interim IRRs early in a fund's life. Ignoring that leverage raises the risk of buyout equity.

\qffield{See also} \hyperref[meth:private-equity-strategies]{strategies over the life cycle}; \hyperref[def:private-debt-features]{private debt}; \hyperref[meth:private-capital-diversification]{diversification}.

\end{qfdefinition}

\begin{qfmethod}[The private investment process: favourable target characteristics, sources of value creation, due diligence and the business plan, and exit routes and their effect on value]\label{meth:private-deal-process-exits}
\qffield{Statement} \emph{Sourcing and selection.} Surveyed buyout firms investigate fewer than 24 of every 100 opportunities in depth, sign fewer than 14 and close 6. They consider about 48 percent of their closed deals proprietary in some way. They rank the business model and competitive position first among selection criteria, then management, their ability to add value, and valuation \citep[NBER WP 21133, pp.~25--26]{w5-gompers-kaplan-mukharlyamov}. Favourable buyout targets have stable cash flows to service debt, assets carried below replacement cost, non-core status for a seller, cost or management improvements available, or owners seeking an exit \citep[pp.~188--189]{ippolito-private-capital-2020}.

\emph{Value creation} (Gompers et al., pp.~27--29; Ippolito, pp.~218--219): revenue growth, which managers expect in about 70 percent of deals; follow-on acquisitions (over 50 percent); cost reduction (36 percent); changes of strategy or management; incentives and governance; buying well; and a well-prepared exit. Kaplan and Str\"omberg's financial, governance and operational engineering summarises the levers \citep[NBER WP 14207, pp.~13--16]{pe-kaplan-stromberg-2009}.

\emph{Due diligence and the business plan.} Business due diligence assesses country, industry and company risks; accounting due diligence assesses quality of earnings, ratios and cash flow; the business plan projects the statements that set the price, the debt and the target return \citep[pp.~71--94, 125--130]{ippolito-private-capital-2020}. Buyout firms value deals mainly with IRR and multiples, targeting a median IRR of about 22--25 percent and a multiple of 2.5, and set exit values from comparable multiples rather than DCF (Gompers et al., pp.~3, 13--17).

\emph{Exit routes} \citep[pp.~220--225]{ippolito-private-capital-2020}:
\begin{itemize}
\item \emph{IPO}: a possibly high price and a staged exit, but costly, and lock-ups leave the fund exposed to the post-IPO price;
\item \emph{trade sale}: a strategic buyer may pay for synergies, with a full cash exit, but sellers resist giving warranties;
\item \emph{secondary buyout}: quicker liquidity and re-leverage when other routes are closed, but with financially minded buyers;
\item \emph{buybacks, options and recapitalisations}: typically used when other routes fail, and usually at low prices.
\end{itemize}
Of exited buyouts in Kaplan and Str\"omberg's sample, 38 percent went to strategic buyers, 24 percent to other funds and 14 percent through IPOs (p.~11).

\qffield{Practice} \emph{Corran Industrial} (original; \hyperref[meth:private-equity-valuation]{valuation}). A carve-out of a non-core division with EBITDA of 50 offers buying below replacement cost, pricing and purchasing gains and two bolt-on acquisitions. Due diligence finds 2 of non-recurring costs and a customer concentration of 30 percent. The business plan assumes 6 percent EBITDA growth.
\begin{itemize}
\item At a 9.0 times EBITDA entry and exit, the equity multiple is 2.03 and the IRR 15.2 percent, below a 22 percent target.
\item Exiting by trade sale at 9.0 times rather than by secondary buyout at 8.0 times is worth 0.32 of multiple, which is why the exit route belongs in the plan from the start.
\end{itemize}

\qffield{Conditions} The survey figures describe what managers say they do, for a 2012--2014 sample; they are not measured outcomes.

\qffield{Pitfalls} Paying for synergies in the entry price that only a strategic buyer can realise. Assuming the exit multiple will exceed the entry multiple. Treating a deal as proprietary when it came through an auction.

\qffield{See also} \hyperref[meth:private-equity-valuation]{valuation}; \hyperref[meth:gp-alignment-net-returns]{GP alignment}.

\end{qfmethod}

\begin{qfmethod}[Private equity strategies over the company life cycle: venture capital, growth equity and buyouts, their characteristics, risks and roles]\label{meth:private-equity-strategies}
\qffield{Statement} Private equity finances companies at each stage, with the instrument and control matched to the risk \citep[pp.~58--59, 111--118]{demaria-private-markets-2021}:
\begin{itemize}
\item \emph{Venture capital} finances companies from seed and early stage through industrialisation to late stage. Minority stakes are bought by capital increases, usually as convertible preferred shares, staged by milestones so that investors can stop funding failures (pp.~112--113). Little or no debt is used; outcomes are highly skewed. Cambridge Associates data give a loss ratio of 35.4 percent for US venture deals (p.~116), and top-5-percent funds an IRR of 64.9 percent against $-16.3$ percent for the bottom 5 percent (p.~113).
\item \emph{Growth equity} finances profitable companies making a transforming move, such as an acquisition, internationalisation or a new product, through minority stakes with little debt. It has a lower loss ratio (13.4 percent) and steadier returns, but often pays rich prices (pp.~114--116).
\item \emph{Buyouts} finance a transfer of ownership of a mature company with equity and debt, through a holding company whose borrowing is serviced by the target. Financial leverage helps when the business earns more than the debt costs and can cost the owner the company when it cannot. Tax deductibility adds tax leverage. Sponsored ownership lasts three to five years (pp.~117--118). The loss ratio is 15.1 percent (p.~116).
\item Turnaround, replacement capital and distressed strategies address companies in difficulty or shareholders wanting a partial exit (pp.~126--127).
\end{itemize}

\emph{Deal terms.} Venture investors use preferred stock, staging, board seats, liquidation preferences and anti-dilution protection \citep[pp.~33--34]{metrick-yasuda-survey-2011}; \citep[pp.~167--173]{ippolito-private-capital-2020}. Buyout investors take control and use leverage, which Sahlman saw as playing a role similar to staging: it removes discretionary cash flow (Metrick and Yasuda, p.~33).

\qffield{Practice} \emph{Fit to strategy} (original). An endowment seeking technology growth and willing to hold for 12 years with manager access can hold venture. An insurer that must limit loss and volatility of capital is better in growth equity and buyouts with modest leverage. A pension worried about drawdowns can add turnaround and distressed managers, whose opportunities rise in recessions.

\qffield{Conditions} Loss ratios and dispersion are Cambridge Associates statistics reported by Demaria and co-authors, as of 2012 and 2016; they depend on vintage and data coverage.

\qffield{Pitfalls} Calling all private equity ``buyouts''. Valuing a venture investment on its last-round price without its preferences. Assuming growth equity is lower risk because it uses less debt, when entry prices are high.

\qffield{See also} \hyperref[def:private-equity-features]{private equity}; \hyperref[meth:private-equity-valuation]{valuation}; \hyperref[meth:special-situations-distressed]{special situations}.

\end{qfmethod}

\begin{qfmethod}[Valuing private equity investments: the venture capital method with dilution, growth equity, and buyout valuation from target returns, leverage and exit multiples]\label{meth:private-equity-valuation}
\qffield{Statement} \emph{Practice of valuation.} Private equity investors value deals mainly by target IRRs and multiples of invested capital, not discounted cash flow; buyout managers report a median target IRR of about 22--25 percent and a multiple of 2.5, and set exit values from comparable multiples \citep[NBER WP 21133, pp.~3, 13--17]{w5-gompers-kaplan-mukharlyamov}.

\emph{Venture capital (and growth capital) method} \citep[pp.~173--175]{ippolito-private-capital-2020}. With investment $I$, target return $r$, holding period $t$ and exit equity value $V$:
\[ R=I(1+r)^t,\qquad S=R/V,\qquad \mathrm{POST}=I/S,\qquad \mathrm{PRE}=\mathrm{POST}-I. \]
New shares for $x$ existing shares are $y=x\,S/(1-S)$, at price $I/y$ (pre-money over existing shares). If later rounds will sell fractions $d_1,d_2$ of the company, the stake bought now must be $S/[(1-d_1)(1-d_2)]$ (Derivation). Targets fall with the stage: above 40 percent for start-ups and above 25 percent for ordinary equity (pp.~20--21).

\emph{Buyouts.} The sponsor works back from the debt the cash flows can carry, the exit multiple and the target return to the price it can pay. Entry and exit multiples are usually assumed equal, and returns are shown across exit years and multiples (pp.~209--210). Equity value creation splits into EBITDA growth, multiple change and debt paydown \citep[pp.~162--163]{demaria-private-markets-2021}.

\qffield{Practice} \emph{Novaro, venture} (original). Invest 8 at a 45 percent target for five years, with exit equity value 200:
\begin{itemize}
\item required value $8\times1.45^5=51.28$, so a 25.64 percent stake at exit;
\item later rounds selling 20 and 15 percent leave 68 percent of today's stake, so 37.70 percent is needed now: post-money 21.22, pre-money 13.22;
\item with 1.0 million existing shares, 0.605 million new shares at 13.22 each.
\end{itemize}
Ippolito adjusts the stake by dividing by $(1+D_1)(1+D_2)$ while describing $D$ as the share taken up by later investors. That gives 18.58 percent, half what is needed under that description; the formula is right only if $D$ measures new shares relative to existing ones.

\emph{Growth equity.} The same method at a 25 percent target and lower dilution gives a much smaller required stake, usually taken as convertible preferred shares with protective rights.

\emph{Corran Industrial, buyout} (original). EBITDA 50, bought at 9.0 times (450) with fees of 10 and 250 of unitranche debt at 9 percent; equity 210. EBITDA grows 6 percent a year, half of it is free cash flow, and all spare cash repays debt:
\begin{center}\small
\begin{tabular}{@{}lrr@{}}
\toprule
 & Exit at 9.0x & Exit at 8.0x \\
\midrule
Exit EBITDA & 66.91 & 66.91 \\
Debt at exit & 176.90 & 176.90 \\
Equity at exit & 425.30 & 358.39 \\
Multiple / IRR & 2.03 / 15.2\% & 1.71 / 11.3\% \\
\bottomrule
\end{tabular}
\end{center}
The 215.30 gain is EBITDA growth times the multiple (152.20) plus debt repaid (73.10) less fees (10). The deal reaches a 22 percent IRR only if it is bought, and sold, at about 6.9 times EBITDA: that is the sponsor's maximum price on these assumptions.

\qffield{Conditions} The exit value is the critical input and rests on comparables; all numbers are hypothetical. The 45 percent target reflects stage risk and the high failure rate, not an expected return.

\qffield{Pitfalls} Ignoring later dilution. Relying on multiple expansion to meet a target return. Equating a high target IRR with the expected return of a portfolio that includes failures.

\qffield{See also} \hyperref[meth:private-equity-strategies]{strategies}; \hyperref[meth:private-debt-financing-structures]{buyout financing}; \hyperref[meth:private-company-valuation]{private company valuation}.

\end{qfmethod}
